\documentclass[12pt]{article}
\usepackage[utf8]{inputenc}
\usepackage[T1]{fontenc}
\usepackage{hyperref}
\title{Urapmin Language Resources for K-12 Curriculum Development}
\author{Compiled by Claude Sonnet 5 (claude-sonnet-5) for Andr\'as Kornai}
\date{}
\begin{document}
\maketitle
\section{Introduction}
Urapmin (ISO 639-3 urm, Glottocode urap1239) is spoken in Telefomin District, Sandaun Province, Papua New Guinea, by roughly 375 people. Sources group it with Tifal, an Ok language. The community is Baptist Christian.
\section{Primary Educational Resources}
None found.
\section{Linguistic and Reference Materials}
\begin{itemize}
\item Urapmin people (Wikipedia): \url{https://en.wikipedia.org/wiki/Urapmin_people}
\item Tifal language (Wikipedia): \url{https://en.wikipedia.org/wiki/Tifal_language}
\item Endangered Languages Project entry: \url{https://endangeredlanguages.com/lang/10694}
\item Telefol dictionary (ANU), for a neighbouring but different language, not Urapmin: \url{https://openresearch-repository.anu.edu.au/items/e2c324d2-d8c7-455c-98de-1cc204051024}
\end{itemize}
No Urapmin grammar or dictionary was found.
\section{Religious and Community Resources}
\begin{itemize}
\item Joshua Project profile: \url{https://joshuaproject.net/people_groups/15724}
\end{itemize}
No Urapmin scripture was found.
\section{Digital and Audio Resources}
None found.
\section{Curriculum Development Context}
Documentation found in search is minimal. The language is not known to be taught in schools. Any curriculum would need to be built from scratch with the community.
\end{document}